\documentclass[11pt, a4paper]{article}
\usepackage[utf8]{inputenc}
\usepackage{pmboxdraw}
\usepackage{newunicodechar}

% Packages for formatting and functionality
\usepackage[utf8]{inputenc}
\usepackage[T1]{fontenc}
\usepackage{geometry}
\usepackage{hyperref}
\usepackage{amsmath}
\usepackage{amssymb}
\usepackage{enumitem}

% Page geometry setup
\geometry{
    left=25mm,
    right=25mm,
    top=25mm,
    bottom=25mm
}

% Hyperlink setup
\hypersetup{
    colorlinks=true,
    linkcolor=blue,
    filecolor=magenta,      
    urlcolor=blue,
    pdftitle={Simple Wikipedia Search Engine},
}

\title{\textbf{Simple Wikipedia Search Engine}}
\author{
    \textbf{Aniruddha Pandey} \\
    \texttt{apandey10@stevens.edu} \\
    \texttt{20037146}
}
\date{\today}

\begin{document}

\maketitle

\noindent
\textbf{Instructor:} Reza Peyrovian \texttt{<rpeyrovi@stevens.edu>} \\
\textbf{Course:} CS 600-A Advanced Algorithm Design and Implementation
\hrule
\newpage

\section{Project Overview}

This project is a search engine built in Rust, designed to index and search a local copy of the Simple English Wikipedia. It consists of multiple components, starting with a high-performance web crawler.

\section{Project Structure}

\begin{verbatim}
.
├── crawler/
│   ├── Cargo.toml
│   └── src/
│       └── main.rs
├── client/
│   ├── package.json
│   └── src/
│       └── routes/
│           └── +page.svelte
├── data/
│   ├── crawled.lst
│   ├── docs.bin
│   ├── inverted_index.bin
│   ├── lemmatization-en.txt
│   ├── seed.lst
│   ├── stopwords-en.txt
│   ├── trie.bin
│   └── whitelist.txt
├── fixtures/
│   ├── 01_basic_search_input.txt
│   ├── 01_basic_search_output.json
│   └── ...
├── screenshots/
│   ├── 01_basic_search_input.png
│   └── ...
├── indexer/
│   ├── Cargo.toml
│   └── src/
│       └── main.rs
├── normalization/
│   ├── Cargo.toml
│   └── src/
│       └── lib.rs
├── server/
│   ├── Cargo.toml
│   └── src/
│       └── main.rs
├── wikipedia-simple-html-dump/
│   └── ... (HTML files)
├── generate_fixtures.py
└── README.tex
\end{verbatim}

\section{Components}

\subsection{Crawler}

The crawler is responsible for traversing the local Wikipedia HTML files to discover and list all available articles.

\subsubsection*{How it Works}
\begin{enumerate}
    \item \textbf{Initialization}: Reads a seed file (\texttt{data/seed.lst}) to find the starting point (e.g., \texttt{simple/index.html}).
    \item \textbf{Traversal}: It uses a Breadth-First Search (BFS) algorithm to visit pages.
    \item \textbf{Parsing}: For each visited page, it parses the HTML to find \texttt{<a>} tags (links).
    \item \textbf{Filtering}: It strictly filters for \textbf{relative local file paths}, ignoring external links (http/https) and anchors (\#).
    \item \textbf{Output}: It generates a list of all discovered unique file paths.
\end{enumerate}

\subsubsection*{Output}
The crawler produces a file named \texttt{crawled.lst} in the \texttt{data/} directory. This file contains the relative paths of all visited documents, stripped of the \texttt{wikipedia-simple-html-dump/} prefix to ensure portability.

Example output in \texttt{data/crawled.lst}:
\begin{verbatim}
simple/index.html
simple/articles/a/r/t/Article_Name.html
...
\end{verbatim}

\subsection{Normalization}
The \texttt{normalization} component is a text processing library responsible for preparing raw text for indexing.

\subsubsection*{Pipeline}
\begin{enumerate}
    \item \textbf{Casefolding}: Converts all text to lowercase to ensure case-insensitive matching.
    \item \textbf{Whitelist Check}: Checks tokens against a whitelist (\texttt{data/whitelist.txt}) to preserve special terms (e.g., "C++", ".NET") that would otherwise be split by the tokenizer.
    \item \textbf{Tokenization}: Splits text into individual tokens (words) based on punctuation and whitespace.
    \item \textbf{Stopword Removal}: Filters out common English filler words (e.g., "the", "is", "at") to reduce index size and improve relevance. Uses \texttt{data/stopwords-en.txt}.
    \item \textbf{Lemmatization}: Reduces words to their root form (e.g., "running" $\rightarrow$ "run") using a dictionary-based approach sourced from \texttt{data/lemmatization-en.txt} (\href{https://github.com/michmech/lemmatization-lists}{Source}).
\end{enumerate}

\subsection{Indexer}
The indexer builds an efficient inverted index from the processed documents.

\subsubsection*{Architecture}
Following the standard inverted index design:
\begin{enumerate}
    \item \textbf{Inverted Index (Occurrence Lists)}: A binary file (\texttt{data/inverted\_index.bin}) storing the list of postings for each term. Each posting contains the Document ID and the list of positions where the term appears (for window scoring).
    \item \textbf{Document Metadata}: A binary file (\texttt{data/docs.bin}) storing metadata for each document, including its path, title, and length (for BM25 scoring).
    \item \textbf{Compressed Trie}: An in-memory compressed trie (implemented using \texttt{trie-rs}) that maps every term in the vocabulary to the index of its occurrence list in the binary file. The trie is serialized to \texttt{data/trie.bin}.
\end{enumerate}

\subsubsection*{Process}
\begin{enumerate}
    \item \textbf{Initialization}: Loads the \texttt{crawled.lst} to map Document IDs (line numbers) to file paths.
    \item \textbf{Processing}: Iterates through every document, parses the HTML to extract text and title, and uses the \texttt{Normalizer} library to tokenize and stem the content.
    \item \textbf{Index Construction}: Builds a temporary in-memory map of \texttt{Term $\rightarrow$ DocID $\rightarrow$ [Positions]}.
    \item \textbf{Optimization}: Sorts postings by Document ID.
    \item \textbf{Serialization}: Saves the occurrence lists, document metadata, and the Trie to disk using \texttt{bincode}.
\end{enumerate}

\subsection{Search \& Ranking}
The search component provides the user interface for querying the index and ranking results.

\subsubsection*{Boolean Logic}
The search engine enforces a strict Boolean AND logic for multi-term queries. Given a query $Q = \{q_1, q_2, \dots, q_n\}$, the result set $R$ is defined as the intersection of the posting lists for each term:
\[ R = \bigcap_{i=1}^{n} \text{Postings}(q_i) \]
This ensures that every returned document contains all the terms in the query. The intersection is computed efficiently by sorting the posting lists by Document ID and using a linear scan algorithm.

\subsubsection*{Ranking Algorithm}
The search engine uses a sophisticated ranking function that combines two relevance signals:

\begin{enumerate}
    \item \textbf{BM25 Score}: A probabilistic information retrieval model that ranks documents based on the query terms appearing in each document.
    
    The BM25 score for a document $D$ and query $Q$ is calculated as:
    \[ \text{BM25}(D, Q) = \sum_{i=1}^{n} \text{IDF}(q_i) \cdot \frac{f(q_i, D) \cdot (k_1 + 1)}{f(q_i, D) + k_1 \cdot (1 - b + b \cdot \frac{|D|}{\text{avgdl}})} \]
    
    Where the Inverse Document Frequency (IDF) is:
    \[ \text{IDF}(q_i) = \ln \left( \frac{N - n(q_i) + 0.5}{n(q_i) + 0.5} + 1 \right) \]
    
    \textbf{Parameters:}
    \begin{itemize}
        \item $f(q_i, D)$: Term frequency of $q_i$ in $D$.
        \item $|D|$: Length of document $D$.
        \item $\text{avgdl}$: Average document length in the corpus.
        \item $N$: Total number of documents.
        \item $n(q_i)$: Number of documents containing $q_i$.
        \item $k_1 = 1.2$, $b = 0.75$: Tunable parameters.
    \end{itemize}

    \item \textbf{Window Score}: Measures the proximity of query terms within the document.
    \[ \text{Window}(D, Q) = \frac{|Q|}{\text{min\_window}(Q, D)} \]
    Where $\text{min\_window}(Q, D)$ is the size of the smallest span of text in $D$ containing all terms in $Q$.
\end{enumerate}

\textbf{Final Score Formula}:
\[ \text{Score}(D, Q) = \alpha \cdot \text{Window}(D, Q) + \beta \cdot \text{BM25}(D, Q) \]
Where $\alpha$ and $\beta$ are weights for the respective scores.

\subsubsection*{Features}
\begin{itemize}
    \item \textbf{Ranked Results}: Returns results sorted by relevance score.
    \item \textbf{Fast Retrieval}: Uses optimized data structures for sub-millisecond query times.
\end{itemize}

\subsection{Server (REST API)}
The server component exposes the search functionality via a RESTful API using \texttt{actix-web}.

\subsubsection*{Endpoints}
\begin{itemize}
    \item \texttt{GET /search?q=<query>}: Performs a search and returns JSON results.
\end{itemize}

\subsubsection*{Response Format}
\begin{verbatim}
{
  "query": "wikipedia",
  "results": [
    { "id": 1, "path": "simple/index.html" },
    ...
  ],
  "total_results": 9145,
  "time_taken_ms": 2.34
}
\end{verbatim}

\subsection{Client (Web UI)}
The client is a modern web interface built with \textbf{SvelteKit} that consumes the REST API.

\subsubsection*{Features}
\begin{itemize}
    \item \textbf{Clean UI}: Simple search bar and results display.
    \item \textbf{Real-time Feedback}: Shows search time and total results found.
    \item \textbf{Direct Links}: Links directly to the local file paths (requires browser permission or local setup).
\end{itemize}

\section{Usage}

\subsection{Prerequisites}
\begin{itemize}
    \item Rust toolchain (Cargo, rustc)
    \item Local Wikipedia dump extracted to \texttt{wikipedia-simple-html-dump/} directory
    \item Lemmatization dictionary at \texttt{data/lemmatization-en.txt} (\href{https://github.com/michmech/lemmatization-lists}{Source})
    \item Stopwords list at \texttt{data/stopwords-en.txt}
    \item Whitelist at \texttt{data/whitelist.txt}
\end{itemize}

\subsection{Configuration}
The crawler uses a \texttt{seed.lst} file in the \texttt{data/} directory to determine the entry point.

\begin{verbatim}
simple/index.html
\end{verbatim}

\subsection{Running the Crawler}
To execute the crawler from the project root. This will traverse the \texttt{wikipedia-simple-html-dump} directory and generate \texttt{data/crawled.lst}.

\begin{verbatim}
cargo run --manifest-path crawler/Cargo.toml
\end{verbatim}

\subsection{Running the Indexer}
To build the index from the crawled data. This will process the files listed in \texttt{crawled.lst} and generate \texttt{data/inverted\_index.bin} and \texttt{data/trie.bin}.

\begin{verbatim}
cargo run --manifest-path indexer/Cargo.toml
\end{verbatim}

\subsection{Running the Server}
To start the REST API server on \texttt{http://127.0.0.1:8080}.

\begin{verbatim}
cargo run --manifest-path server/Cargo.toml
\end{verbatim}

\subsection{Running the Client}
To start the web interface. Ensure you have Node.js installed.

\textbf{Important:} The client requires a symbolic link to the Wikipedia dump to serve the files locally.

\begin{enumerate}
    \item Create the symbolic link (if not already created):
    \begin{verbatim}
# Run from the project root
ln -s "../../wikipedia-simple-html-dump" "client/static/wiki"
    \end{verbatim}
    
    \item Install dependencies and run the development server:
    \begin{verbatim}
cd client
npm install
npm run dev
    \end{verbatim}
\end{enumerate}

The client will be available at \texttt{http://localhost:5173}. Clicking on search results will open the local Wikipedia pages directly in your browser.

\subsection{Running Tests}
To run the unit tests for the Normalizer library (lemmatization, tokenizer, etc.):

\begin{verbatim}
cargo test --manifest-path normalization/Cargo.toml
\end{verbatim}

\subsection{Generating Test Fixtures}
A Python script is included to generate a set of test cases for the search API, covering various boundary conditions. This script requires the server to be running.

\begin{enumerate}
    \item Ensure the server is running (see above).
    \item Run the script from the project root:
    \begin{verbatim}
python3 generate_fixtures.py
    \end{verbatim}
\end{enumerate}

This will create a `fixtures/` directory populated with input (`.txt`) and output (`.json`) files for each test case, such as basic searches, empty queries, and queries with special characters.

\subsubsection*{Test Case Screenshots}
The `screenshots/` directory contains captured images of these test cases, demonstrating the system's correct handling of various boundary conditions and edge cases.

\section{Data Source}

The data used for this search engine is sourced from the Wikimedia Dumps, specifically the Simple English Wikipedia Database dump, which can be found at \href{https://dumps.wikimedia.org/}{dumps.wikimedia.org}.
The lemmatization list is sourced from \href{https://github.com/michmech/lemmatization-lists}{michmech/lemmatization-lists}.
The stopwords list is sourced from the `nltk` python library.

\newpage

\section*{Academic Integrity Pledge}
\emph{I pledge my honor that I have abided by the Stevens Honor System.}\footnote{\href{https://library.stevens.edu/academicintegrity\#s-lg-link-list-70154023}{library.stevens.edu/academicintegrity}}

\end{document}